\subsubsection{Bellman-Ford}
\begin{lstlisting}[style=mystyle, language=C++]
// TC: O(N * M)
// 1. Camino minimo desde s y deteccion de ciclo negativo alcanzable
vector<long long> bellman_ford(int n, vector<array<int,3>>& edges, int s, bool& has_cycle){
	const long long INF = 1e18;
	vector<long long> dist(n, INF);
	dist[s] = 0;
	for(int i = 0; i < n - 1; i++){
		for(auto [u, v, w] : edges){
			if(dist[u] < INF && dist[u] + w < dist[v]){
				dist[v] = max(-INF, dist[u] + w);
			}
		}
	}
	has_cycle = false;
	for(auto [u, v, w] : edges){
		if(dist[u] < INF && dist[u] + w < dist[v]){
			has_cycle = true; break;
		}
	}
	return dist;
}

// 2. Reconstruir ciclo negativo (ej. CSES - Cycle Finding)
// Retorna vertices del ciclo {v1, v2, ..., v1} en orden dirigido, o vacio si no hay.
// s = -1: detecta CUALQUIER ciclo en el grafo (dist=0 para todo nodo).
// s >= 0: detecta ciclo negativo alcanzable desde s.
vector<int> get_negative_cycle(int n, vector<array<int,3>>& edges, int s = -1){
	const long long INF = 1e18;
	vector<long long> dist(n, s == -1 ? 0 : INF);
	vector<int> parent(n, -1);
	if(s != -1) dist[s] = 0;
	int last_relaxed = -1;

	// N relajaciones: si en la vuelta N se relaja algo, hay ciclo negativo
	for(int rep = 0; rep < n; rep++){
		last_relaxed = -1;
		for(auto [u, v, w] : edges){
			if(dist[u] < INF && dist[u] + w < dist[v]){
				dist[v] = max(-INF, dist[u] + w);
				parent[v] = u;
				last_relaxed = v;
			}
		}
	}
	if(last_relaxed == -1) return {}; // No hay ciclo negativo

	// Retroceder N veces: el ultimo vertice relajado podria estar fuera del ciclo
	// (en una rama alcanzable). Al retroceder N veces por parent se entra seguro al ciclo.
	for(int i = 0; i < n; i++) last_relaxed = parent[last_relaxed];

	vector<int> cycle;
	for(int cur = last_relaxed; ; cur = parent[cur]){
		cycle.push_back(cur);
		if(cur == last_relaxed && cycle.size() > 1) break;
	}
	reverse(cycle.begin(), cycle.end());
	return cycle;
}
\end{lstlisting}
